\section{Point targets and implementation extensions}
\label{sec:supp-extensions}

\subsection{Localization bias and conditional quantiles}

\begin{lemma}[Uniform localization bias]
\label{lem:bias}
Under Conditions~\ref{cond:kernel} and \ref{cond:bias},
\[
\sup_{y\in\mathcal Y}
|F_{h,0}(y\mid x_0)-F_0(y\mid x_0)|
=O(h^2).
\]
\end{lemma}

\begin{proof}
Let \(g_y(x)=q_{0,y}(x)f_X(x)\). Then
\[
E(K_hZ_y)
=h^d\int K(u)g_y(x_0+hu)\,du,
\qquad
\mu_h=h^d\int K(u)f_X(x_0+hu)\,du.
\]
Uniform second-order Taylor expansions and symmetry of \(K\) give
\begin{align*}
E(K_hZ_y)
&=h^d\{\kappa_1g_y(x_0)+O(h^2)\},\\
\mu_h
&=h^d\{\kappa_1f_X(x_0)+O(h^2)\},
\end{align*}
uniformly in \(y\). Taking the ratio proves the result.
\end{proof}

\begin{proof}[Proof of the point-target statement in Corollary~\ref{cor:F0}]
By Lemma~\ref{lem:bias},
\[
\sqrt{N_{M,h}}
\|F_{h,0}-F_0\|_{\mathcal Y}
=O\{\sqrt{N_{M,h}}h^2\}=o(1).
\]
The process and multiplier conclusions follow from Theorems~\ref{thm:process} and \ref{thm:bootstrap} and Slutsky's theorem in \(\ell^\infty(\mathcal Y)\).
\end{proof}

\begin{proposition}[Increasing rearrangement]
\label{prop:rearrangement}
Suppose \(F_0(\cdot\mid x_0)\) is continuously differentiable and strictly increasing on the threshold region, with derivative bounded away from zero and values in \([\eta,1-\eta]\). If
\[
r_M(\widehat F_h^{\mathrm{raw}}-F_0)
\Rightarrow\mathbb G_p
\quad\text{in }\ell^\infty(\mathcal Y),
\qquad
r_M=\sqrt{N_{M,h}},
\]
then
\begin{equation}
r_M(\widehat F_h^{\mathrm{mon}}-F_0)
=
r_M(\widehat F_h^{\mathrm{raw}}-F_0)
+o_{\mathbb P}(1)
\quad\text{in }\ell^\infty(\mathcal Y),
\label{eq:rearrangement-equivalence}
\end{equation}
and therefore
\[
r_M(\widehat F_h^{\mathrm{mon}}-F_0)
\Rightarrow\mathbb G_p
\quad\text{in }\ell^\infty(\mathcal Y).
\]
For \(1\le r\le\infty\), rearrangement does not increase the \(L_r(\nu)\) distance to the monotone target.
\end{proposition}

\begin{proof}
The rearrangement contraction and functional limit result at a strictly increasing target follow from \citet{ChernozhukovEtAl2010}. At such a target, the Hadamard derivative of increasing rearrangement is the identity. Clipping is inactive with probability tending to one on the stated interior probability region. The functional delta method gives \eqref{eq:rearrangement-equivalence} and the asserted weak limit.
\end{proof}

\begin{proof}[Proof of the conditional-quantile statement in Corollary~\ref{cor:F0}]
Define the point-target empirical processes
\[
\mathbb Z_{M,0}^{\mathrm{raw}}
=
\sqrt{N_{M,h}}(\widehat F_h^{\mathrm{raw}}-F_0),
\qquad
\mathbb Z_{M,0}^{\mathrm{mon}}
=
\sqrt{N_{M,h}}(\widehat F_h^{\mathrm{mon}}-F_0).
\]
By \eqref{eq:rearrangement-equivalence},
\[
\|\mathbb Z_{M,0}^{\mathrm{mon}}
 -\mathbb Z_{M,0}^{\mathrm{raw}}\|_{\mathcal Y}
=o_{\mathbb P}(1).
\]
The point-target statement proved above gives
\[
\mathbb Z_{M,0}^{\mathrm{raw}}
\Rightarrow\mathbb G_p
\quad\text{in }\ell^\infty(\mathcal Y).
\]
The inverse-CDF map is Hadamard differentiable at \(F_0\), tangentially to continuous functions, with derivative
\[
g\mapsto
-
\frac{g\{Q_0(\tau\mid x_0)\}}
{f_{Y\mid X}\{Q_0(\tau\mid x_0)\mid x_0\}}.
\]
The functional delta method, Theorem 3.9.4 of \citet{VanderVaartWellner1996}, yields the uniform Bahadur representation
\begin{equation}
\sup_{\tau\in\mathcal T}
\left|
\sqrt{N_{M,h}}\{\widehat Q_h(\tau\mid x_0)-Q_0(\tau\mid x_0)\}
+
\frac{\mathbb Z_{M,0}^{\mathrm{raw}}\{Q_0(\tau\mid x_0)\}}
{f_{Y\mid X}\{Q_0(\tau\mid x_0)\mid x_0\}}
\right|
=o_{\mathbb P}(1).
\label{eq:uniform-bahadur}
\end{equation}
Combining \eqref{eq:uniform-bahadur} with the point-target process limit proves \eqref{eq:quantile-limit}--\eqref{eq:quantile-gaussian-limit}.
\end{proof}

\subsection{Estimated verification probabilities}

\begin{lemma}[Exact propensity drift]
\label{lem:propensity-drift}
Let \(\widehat\pi\) be fitted independently of the observation on which \(U_y(a,\widehat\pi)\) is evaluated. Then
\begin{equation}
E\{U_y(a,\widehat\pi)-m_{0,y}\mid W,\widehat\pi\}
=
(a_y-m_{0,y})
\frac{\widehat\pi-\pi_M}{\widehat\pi}.
\label{eq:supp-pi-drift}
\end{equation}
\end{lemma}

\begin{proof}
Conditional on \(W\) and \(\widehat\pi\),
\[
\begin{aligned}
E\{U_y(a,\widehat\pi)\mid W,\widehat\pi\}
&=a_y+\frac1{\widehat\pi}E[R(Z_y-a_y)\mid W]\\
&=a_y+\frac{\pi_M}{\widehat\pi}(m_{0,y}-a_y).
\end{aligned}
\]
Subtracting \(m_{0,y}\) gives \eqref{eq:supp-pi-drift}.
\end{proof}

\begin{proposition}[Estimated-propensity process extension]
\label{prop:pi-extension}
Suppose the conditions of Theorem~\ref{thm:process} hold for the estimator using \(\pi_M\). Let \(\widehat\pi\) be fitted on a separate fold and satisfy, uniformly over \(y\in\mathcal Y\),
\begin{align}
&\sqrt{N_{M,h}}
\left|
\frac1{\mu_h}
E\left[
K_h(a_y-m_{0,y})
\frac{\widehat\pi-\pi_M}{\widehat\pi}
\,\Bigm|\,\mathcal T
\right]
\right|
=o_{\mathbb P}(1),
\label{eq:pi-drift-cond}\\
&\frac{p_Mh^d}{\mu_h^2}
E\left[
K_h^2R(Z_y-a_y)^2
\left(\frac1{\widehat\pi}-\frac1{\pi_M}\right)^2
\,\Bigm|\,\mathcal T
\right]
=o_{\mathbb P}(1),
\label{eq:pi-var-cond}
\end{align}
and suppose the corresponding random class satisfies Condition~\ref{cond:entropy}. Then replacing \(\pi_M\) by \(\widehat\pi\) leaves the limit in Theorem~\ref{thm:process} unchanged.
\end{proposition}

\begin{proof}
Decompose
\[
U_y(a,\widehat\pi)-U_y(a,\pi_M)
=
R(Z_y-a_y)
\left(\frac1{\widehat\pi}-\frac1{\pi_M}\right).
\]
Lemma~\ref{lem:propensity-drift} identifies the conditional mean. Condition \eqref{eq:pi-drift-cond} removes the localized mean at the process scale. After centering, Condition \eqref{eq:pi-var-cond} and the entropy bound make the empirical process \(o_{\mathbb P}(1)\) by the maximal-inequality argument in Lemma~\ref{lem:nuisance-perturb}. The process expansion using \(\pi_M\) therefore remains unchanged.
\end{proof}

Define
\begin{align*}
r_{a,M,h}^2
&=
\sup_y
\frac1{\mu_h}
E\{K_h(a_y-m_{0,y})^2\mid\mathcal T\},\\
r_{\pi,M,h}^2
&=
\frac1{\mu_h}
E\left[
K_h
\left(\frac{\widehat\pi-\pi_M}{\widehat\pi}\right)^2
\Bigm|\mathcal T
\right].
\end{align*}
Cauchy--Schwarz gives
\[
\sup_y
\left|
\frac1{\mu_h}
E\left[
K_h(a_y-m_{0,y})
\frac{\widehat\pi-\pi_M}{\widehat\pi}
\Bigm|\mathcal T
\right]
\right|
\le r_{a,M,h}r_{\pi,M,h}.
\]
Thus
\[
\sqrt{N_{M,h}}\,r_{a,M,h}r_{\pi,M,h}
=o_{\mathbb P}(1)
\]
is sufficient for \eqref{eq:pi-drift-cond}. Along an interior limiting path,
\[
a_{\lambda,y}-m_{0,y}
\to
(1-\lambda)\{q_{0,y}-m_{0,y}\},
\]
so the relevant propensity rate is a localized relative-error rate.

\subsection{Grid approximation}

Let \(\nu\) have a bounded density on \(\mathcal Y\), let
\[
\underline y=y_{1,L}<\cdots<y_{L,L}=\overline y,
\]
and let \(\nu_L=\sum_{j=1}^Lw_{j,L}\delta_{y_{j,L}}\) be a positive quadrature rule with mesh \(\Delta_L\).

\begin{proposition}[Grid-to-continuum transfer]
\label{prop:quadrature}
Under Conditions~\ref{cond:sampling} and \ref{cond:kernel}, suppose, uniformly
in the local region,
\[
|A_yD_y-A_zD_z|+|D_y^2-D_z^2|
\le C_\beta|y-z|^\beta
\]
for some \(\beta>0\) and \(C_\beta<\infty\). Suppose the quadrature rule has order \(\beta\). Let \((G_{M,h},H_{M,h})\) use \(\nu\) and \((G_{M,h,L},H_{M,h,L})\) use \(\nu_L\). Then
\[
|G_{M,h,L}-G_{M,h}|+|H_{M,h,L}-H_{M,h}|
\le
C\Delta_L^\beta E v_{M,h}.
\]
After multiplication by \(p_M\), the right side is \(O(\Delta_L^\beta)\). If \(\inf_M\mathcal H_{M,h}>0\),
\[
|\lambda_{h,L}^{\mathrm{or}}-\lambda_h^{\mathrm{or}}|
=O(\Delta_L^\beta).
\]
\end{proposition}

\begin{proof}
The quadrature error for a uniformly \(\beta\)-H\"older function is \(O(\Delta_L^\beta)\). Apply this result pointwise to \(A_yD_y\) and \(D_y^2\), multiply by \(v_{M,h}\), and take expectation. Since \(p_ME v_{M,h}=O(1)\), the normalized error is \(O(\Delta_L^\beta)\). Projection is nonexpansive and the ratio map is Lipschitz when \(\mathcal H\) is bounded below.
\end{proof}

\begin{proposition}[Growing-grid approximation of process suprema]
\label{prop:growing-grid}
Let \(\mathcal Y_{L_M}\subset\mathcal Y\) be deterministic grids with mesh \(\Delta_{L_M}\to0\). Under Theorems~\ref{thm:process} and \ref{thm:bootstrap},
\begin{align*}
\left|
\|\mathbb Z_M\|_{\mathcal Y}
-
\max_{y\in\mathcal Y_{L_M}}|\mathbb Z_M(y)|
\right|&=o_{\mathbb P}(1),\\
\left|
\|\mathbb G_M^\xi\|_{\mathcal Y}
-
\max_{y\in\mathcal Y_{L_M}}|\mathbb G_M^\xi(y)|
\right|&=o_{\mathbb P}(1)
\end{align*}
conditionally in probability for the multiplier process. If the reported CDF is extended by piecewise-linear interpolation and \(F_{h,0}\) has a uniformly bounded derivative, the deterministic interpolation error is negligible whenever
\[
\sqrt{N_{M,h}}\,\Delta_{L_M}\to0.
\]
\end{proposition}

\begin{proof}
The process and multiplier proofs establish asymptotic equicontinuity under the totally bounded semimetric \(\rho\). The maximal \(\rho\)-distance from any point in \(\mathcal Y\) to the grid tends to zero, and each supremum difference is bounded by the corresponding modulus of continuity. Piecewise-linear interpolation of the target has error at most \(C\Delta_{L_M}\), while stochastic interpolation is controlled by the same process modulus.
\end{proof}

\subsection{Local-linear equivalent-kernel extension}

The local-constant construction extends to local-polynomial distribution estimation and related boundary-adaptive conditional-density procedures through equivalent-kernel representations; see \citet{CattaneoJanssonMa2024} and \citet{CattaneoEtAl2024}. We give the interior local-linear version for the selective-verification model.

Let
\[
r_h(X)=r\{(X-x_0)/h\},
\qquad
r(u)=(1,u^\top)^\top,
\]
and define
\begin{align}
\mathbf S_h
&=E\{K_h(X-x_0)r_h(X)r_h(X)^\top\},
\label{eq:ll-gram}\\
\gamma_{h,y}
&=E\{K_h(X-x_0)r_h(X)Z_y\},
\qquad
\beta_{h,y}=\mathbf S_h^{-1}\gamma_{h,y}.
\label{eq:ll-beta}
\end{align}
The local-linear regularized target is
\begin{equation}
F_{h,0}^{\mathrm{LL}}(y\mid x_0)
=e_1^\top\beta_{h,y},
\qquad
e_1=(1,0,\ldots,0)^\top\in\mathbb R^{d+1}.
\label{eq:ll-target}
\end{equation}
For a path \(a\) fixed conditional on data outside the evaluation sample, let
\begin{align*}
\widehat{\mathbf S}_h
&=\mathbb P_M\{K_hr_hr_h^\top\},\\
\widehat\gamma_{h,y}^{a}
&=\mathbb P_M\{K_hr_hU_y(a)\},\\
\widehat\beta_{h,y}^{a}
&=\widehat{\mathbf S}_h^{-1}\widehat\gamma_{h,y}^{a},\\
\widehat F_h^{\mathrm{LL},a}(y\mid x_0)
&=e_1^\top\widehat\beta_{h,y}^{a}.
\end{align*}
The sample-split estimator applies this construction foldwise and combines the evaluation folds as in \eqref{eq:final-estimator}. Define the population equivalent weight
\begin{equation}
L_h(X;x_0)
=e_1^\top\mathbf S_h^{-1}r_h(X)K_h(X-x_0).
\label{eq:ll-weight}
\end{equation}

\begin{proposition}[Local-linear efficiency, transfer, and process]
\label{prop:local-linear}
Under Condition~\ref{cond:sampling}, fix \(M\) and \(h\) such that
\(\mathbf S_h\) is nonsingular and
\[
E\{L_h^2(X;x_0)/\pi_M(W)\}<\infty.
\]
For a finite threshold vector, the canonical-gradient components are
\begin{equation}
\phi_{M,h,y}^{\mathrm{LL},\mathrm{eff}}(O)
=
L_h(X;x_0)
\left[
 m_{0,y}(W)
 +\frac{R}{\pi_M(W)}\{Z_y-m_{0,y}(W)\}
 -r_h(X)^\top\beta_{h,y}
\right].
\label{eq:ll-efficient-if}
\end{equation}
For a generic augmentation, set
\begin{equation}
\phi_{M,h,y}^{\mathrm{LL},a}(O)
=
L_h(X;x_0)
\{U_y(a)-r_h(X)^\top\beta_{h,y}\},
\label{eq:ll-influence}
\end{equation}
and
\[
\Omega_{M,h}^{\mathrm{LL},a}(y,z)
=
h^dE\{\phi_{M,h,y}^{\mathrm{LL},a}
        \phi_{M,h,z}^{\mathrm{LL},a}\}.
\]
Under a known verification law, \(\phi^{\mathrm{LL},a}\) is an influence function and
\begin{equation}
\Omega_{M,h}^{\mathrm{LL},a}(y,z)
-
\Omega_{M,h}^{\mathrm{LL},m_0}(y,z)
=
h^dE\left[
L_h^2(X;x_0)
\left\{\frac1{\pi_M(W)}-1\right\}
\delta_y^a(W)\delta_z^a(W)
\right].
\label{eq:ll-psd}
\end{equation}
The difference is positive semidefinite over every finite threshold collection, and \(\phi^{\mathrm{LL},m_0}=\phi^{\mathrm{LL},\mathrm{eff}}\) remains canonical when the verification law is unrestricted.

Suppose next that Condition~\ref{cond:kernel} holds and that the
sample-separation, curvature, and candidate-path clauses of
Condition~\ref{cond:candidates} hold for the local-linear equivalent-kernel
metric. Suppose that
\[
h^{-d}\mathbf S_h\to f_X(x_0)\mathbf S_K,
\qquad
\mathbf S_K=\int K(u)r(u)r(u)^\top\,du,
\]
where \(\mathbf S_K\) is positive definite. Assume the equivalent-kernel
versions of Conditions~\ref{cond:nuisance}, \ref{cond:entropy},
\ref{cond:path-compatibility}, and \ref{cond:covariance} hold for \(L_h\), the
fitted paths, and their threshold increments. Then
\begin{equation}
\sqrt{N_{M,h}}
\{\widehat F_h^{\mathrm{LL}}(a_M;\cdot\mid x_0)
-F_{h,0}^{\mathrm{LL}}(\cdot\mid x_0)\}
\Rightarrow
\mathbb G_p^{\mathrm{LL}}
\quad\text{in }\ell^\infty(\mathcal Y).
\label{eq:ll-process-limit}
\end{equation}
The transfer criterion, observable gain representation, coefficient estimator, path-oracle equivalence, and multiplier process remain valid after replacing \(v_{M,h}\) by
\begin{equation}
v_{M,h}^{\mathrm{LL}}(W)
=h^dL_h^2(X;x_0)
\left\{\frac1{\pi_M(W)}-1\right\}.
\label{eq:ll-transfer-weight}
\end{equation}
If
\[
\ell_K(u)=e_1^\top\mathbf S_K^{-1}r(u)K(u),
\qquad
\kappa_K^{\mathrm{LL}}=\int\ell_K^2(u)\,du,
\]
then the covariance kernel in Theorem~\ref{thm:process} has \(\kappa_K\) replaced by \(\kappa_K^{\mathrm{LL}}\). Under Condition~\ref{cond:bias},
\[
\sup_{y\in\mathcal Y}
|F_{h,0}^{\mathrm{LL}}(y\mid x_0)-F_0(y\mid x_0)|
=O(h^2),
\]
so the point-target and quantile conclusions in Corollary~\ref{cor:F0} apply to the local-linear estimator.
\end{proposition}

\begin{proof}
The augmentation identity gives
\[
E\{K_hr_hU_y(a)\}=E(K_hr_hZ_y)=\gamma_{h,y},
\]
so \(F_{h,0}^{\mathrm{LL}}\) is independent of the augmentation. The parameter \(\beta_{h,y}\) solves
\[
E[K_hr_h\{Z_y-r_h^\top\beta_{h,y}\}]=0.
\]
Differentiating this moment equation along a full-data parametric submodel with score \(s\) gives
\[
\dot\beta_{h,y}
=
\mathbf S_h^{-1}
E[K_hr_h\{Z_y-r_h^\top\beta_{h,y}\}s].
\]
Hence the full-data gradient for \(F_{h,0}^{\mathrm{LL}}(y\mid x_0)\) is
\[
L_h(X;x_0)\{Z_y-r_h(X)^\top\beta_{h,y}\}.
\]
Applying the observed-data tangent decomposition in Section~\ref{sec:supp-efficiency} replaces this gradient by its conditional mean given \(W\) plus the inverse-probability residual, which is exactly \eqref{eq:ll-efficient-if}. The same propensity-tangent calculation as in Theorem~\ref{thm:signed-efficiency} proves canonicality under an unrestricted verification law.

The sample and population normal equations give the exact identity
\begin{equation}
\widehat\beta_{h,y}^{a}-\beta_{h,y}
=
\widehat{\mathbf S}_h^{-1}
(\mathbb P_M-P)
\left[K_hr_h\{U_y(a)-r_h^\top\beta_{h,y}\}\right].
\label{eq:ll-exact-linearization}
\end{equation}
Every entry of \(K_hr_hr_h^\top\) is bounded and supported on a set of probability order \(h^d\). Therefore
\begin{equation}
\|\widehat{\mathbf S}_h-\mathbf S_h\|
=O_{\mathbb P}\{(h^d/M)^{1/2}\},
\qquad
\|\mathbf S_h^{-1}(\widehat{\mathbf S}_h-\mathbf S_h)\|
=O_{\mathbb P}\{(Mh^d)^{-1/2}\}.
\label{eq:ll-gram-rate}
\end{equation}
The matrix inverse expansion is
\begin{equation}
\widehat{\mathbf S}_h^{-1}
=
\mathbf S_h^{-1}
-
\mathbf S_h^{-1}(\widehat{\mathbf S}_h-\mathbf S_h)\mathbf S_h^{-1}
+\mathcal R_{M,h}^{\mathbf S},
\label{eq:ll-inverse-expansion}
\end{equation}
with
\[
\|\mathcal R_{M,h}^{\mathbf S}\|
=O_{\mathbb P}\{h^{-d}(Mh^d)^{-1}\}.
\]
The empirical score in \eqref{eq:ll-exact-linearization} is
\[
O_{\mathbb P}\{(h^d/(Mp_M))^{1/2}\}
\]
in each fixed direction and uniformly over thresholds under Condition~\ref{cond:entropy}. Thus
\[
\|\widehat{\mathbf S}_h^{-1}-\mathbf S_h^{-1}\|
\,
O_{\mathbb P}\{(h^d/(Mp_M))^{1/2}\}
=
O_{\mathbb P}\left\{
N_{M,h}^{-1/2}(Mh^d)^{-1/2}
\right\}
=o_{\mathbb P}(N_{M,h}^{-1/2}).
\]
Equation \eqref{eq:ll-influence} is therefore the first-order influence function.

The difference between \(\phi^{\mathrm{LL},a}\) and \(\phi^{\mathrm{LL},m_0}\) is
\[
L_h(X;x_0)\delta_y^a(W)
\left(1-\frac R{\pi_M(W)}\right).
\]
It is orthogonal to the known-design tangent space. Lemma~\ref{lem:cov-decomp} gives \eqref{eq:ll-psd}, and its finite-dimensional quadratic forms are nonnegative. Integration over \(\nu\) produces the quadratic transfer criterion with weight \eqref{eq:ll-transfer-weight}; Lemma~\ref{lem:observable-G} gives the observable alignment representation. The gain-estimation, coefficient, path-replacement, and multiplier arguments then apply foldwise with the equivalent weight.

Finally,
\[
h^{-d}\mathbf S_h\to f_X(x_0)\mathbf S_K
\]
implies
\[
L_h(x_0+hu;x_0)
=
h^{-d}f_X(x_0)^{-1}\ell_K(u)+o(h^{-d})
\]
uniformly on the support of \(K\). The change-of-variables argument in Lemma~\ref{lem:covariance-limit} replaces \(\kappa_K\) by \(\kappa_K^{\mathrm{LL}}\). The envelope, increment, and entropy calculations retain the orders established in Sections~\ref{sec:supp-process}--\ref{sec:supp-bootstrap}, proving the process and multiplier limits. A uniform second-order expansion of \(q_{0,y}(x)\) around \(x_0\), combined with the local-linear moment equations, gives the \(O(h^2)\) bias.
\end{proof}
